\section{Conclusion}\label{sec:conclusion}

This paper asks whether the settlement window of the National Stock Exchange of India behaves
differently on expiry days, and whether it shows the marks of trading designed to move the
settlement price. The original thirty hypotheses are consolidated into four, each tested through
one primary measure, with the remaining measures as supporting evidence.

\paragraph{BH1, the settlement price: not supported.} The settlement VWAP ends no further from the
mid at the start of the window on expiry days. It ends further from the last traded price, more
so in securities with derivatives, which describes a more volatile end to the window. The stock
settlement is partly reversed the next morning, but equally in securities without derivatives,
so the reversal is a property of the expiry session rather than of the settlement. The expiring stock
future is priced on the cash price until the close, as settlement by delivery implies, and
settlement prices are not drawn toward option strikes. The exception is the index: the Nifty 50,
whose derivatives are settled in cash, falls over the window on index expiries relative to an
index without derivatives.

\paragraph{BH2, the cost of trading: supported.} The quoted spread widens, displayed depth falls
and a given flow moves the price further, and the widening of the spread is specific to
securities with derivatives.

\paragraph{BH3, directional pressure: not supported.} The absolute order imbalance in the window
does not change. Imbalance persists more through the window on monthly expiries, as the
execution of large orders would produce, but in all groups alike; the variance ratio and the
permanent share of price impact move, relative to the placebo group, toward congestion rather
than pressure.

\paragraph{BH4, order behaviour and participation: supported.} More orders are withdrawn within a
second of entry, specifically in securities with derivatives, while cancellations per entry,
concealment and the use of immediate-or-cancel orders do not change. Institutional participation
in the window rises on monthly expiries, and proprietary volume rises with it.

Taken together, expiry draws more trading into the settlement window and makes it more expensive
to trade in, through a more fleeting order book, without the one-sided order flow that pushing
the settlement price would require. The incentive to push is weak for stock derivatives because
they were settled by delivery, and the one settlement that is displaced is the one paid in cash.

The analysis has limitations. The order book sample covers twelve expiries and their controls;
the derivatives files cover four expiries. Open interest is not observed, so positions cannot be
tied to trading in the window. The participant categories do not identify individual traders, so
a strategy spread across participants is not detected as such. The closing value of the index is
approximated by its average over the window. And a systematic effect in the average is not
evidence of manipulation by any particular trader, which requires surveillance data that
identify the trader.
